\documentclass[10pt]{article}
\usepackage[margin=0.70in]{geometry}
\usepackage[T1]{fontenc}
\usepackage[utf8]{inputenc}
\usepackage[table]{xcolor}
\usepackage{booktabs}
\usepackage{array}
\usepackage{tabularx}
\usepackage[hidelinks]{hyperref}
\setlength{\parindent}{0pt}
\setlength{\parskip}{0.38em}
\setlength{\emergencystretch}{2em}

\newcommand{\questionsep}{\par\vspace{1.25em}}
\definecolor{inhouseblue}{RGB}{225,240,255}

\begin{document}

\begin{center}
{\Large Response to Reviewer 7Gar}\\[2pt]
Submission 13175
\end{center}

\textbf{Opening response.}
We thank the reviewer for identifying where our framing ran ahead of the
evidence. We will reorganize each experiment around its question, unit, input,
target, metric, and conclusion, and narrow claims for derived targets. This
resource has accompanied me since my first PhD year; near graduation, I hope its
unified release helps new students enter spatial transcriptomics.

\questionsep
\textbf{Q1.}\quad
\textit{In-house data scale. Please quantify precisely how much new in-house
data is contributed relative to HEST-1K and STimage-1K4M (samples, tissues,
platforms, spots/cells), ideally as a comparison table. What fraction of the
resource is genuinely new vs.\ aggregated?}

\textbf{A1.}\enspace
We agree and will distinguish primary acquisition from aggregation. We audited
each sample against HEST-1K v1.1.0 and the 2025-02-12 STimage-1K4M snapshot.

\begin{center}
\begingroup
\footnotesize
\setlength{\tabcolsep}{3pt}
\renewcommand{\arraystretch}{1.12}
\begin{tabularx}{\linewidth}{
  >{\raggedright\arraybackslash}p{0.22\linewidth}
  >{\raggedright\arraybackslash}p{0.12\linewidth}
  >{\raggedright\arraybackslash}X
  >{\raggedright\arraybackslash}p{0.15\linewidth}
  >{\raggedleft\arraybackslash}p{0.17\linewidth}}
\toprule
Category & Slides/samples & Organs / cell-line origin & Platforms &
Aligned target cells after QC \\
\midrule
\rowcolor{inhouseblue}
New in-house primary data &
21 &
mouse-embryo; cervical; lung &
DBiC-seq &
53,989 \\
Public, overlapping HEST-1K &
32 unique slides (33 records) &
22 species--organ strata &
Xenium; Visium HD &
11,051,149 \\
Public, overlapping STimage-1K4M only &
2 &
2 species--organ strata &
Visium HD &
213,405 \\
Other public data &
44 &
16 species--organ strata &
Xenium; Visium HD &
16,996,704 \\
\textbf{Total sMMC-28M} &
\textbf{99 unique samples (100 records)} &
\textbf{25 tissue strata + 3 cell-line origins} &
\textbf{Xenium; Visium HD; DBiC-seq} &
\textbf{28,315,247} \\
\bottomrule
\end{tabularx}
\endgroup
\end{center}

\textbf{A2.}\enspace The audit establishes two boundaries:
\begin{itemize}
\item \textbf{New in-house data:} our collaborators provided the DBiC-seq data.
Each released record is a native, individually resolved cell paired one-to-one
with its transcriptomic profile, making this a high-quality source of
cell-aligned supervision.
\item \textbf{Limited overlap:} only 32 examples overlap HEST-1K, all from its
newly added Visium HD cohort; approximately 70\% of our collection is not
contained in HEST-1K.
\end{itemize}



\questionsep
\textbf{Q2.}\quad
\textit{Clear purpose and biological meaning of the per-axis experiments.}

\textbf{A1 (Scale).}\enspace
We agree that the three axes lacked clear objectives and conclusions and will
reorganize them as follows. This experiment is not intended as a benchmark; it
tests the central scale claim of the resource: whether increasing the amount of
training data improves model performance.
\begin{itemize}
\item \textbf{Objective:} determine whether a larger training set from our
dataset benefits image-to-expression models.
\item \textbf{Design/metric:} the original Fig.\ 2 varies the training fraction
while fixing the targets, preprocessing, model head, and spatial test regions,
and evaluates gene-macro Pearson across held-out cells. We further tested seven
frozen encoders with training-selected top-50 HVGs, four buffered spatial
holdouts, and 5\%/10\%/25\%/100\% training fractions.
\end{itemize}

\begin{center}
\begingroup
\scriptsize
\setlength{\tabcolsep}{4pt}
\renewcommand{\arraystretch}{1.05}
\begin{tabular*}{\linewidth}{@{\extracolsep{\fill}}lrrrr@{}}
\toprule
Encoder & 5\% & 10\% & 25\% & 100\% \\
\midrule
ResNet-50 & 0.140 & 0.160 & 0.190 & 0.221 \\
CTransPath & 0.162 & 0.186 & 0.222 & 0.259 \\
Phikon & 0.168 & 0.200 & 0.237 & 0.273 \\
CONCH & 0.178 & 0.213 & 0.256 & 0.292 \\
UNI & 0.188 & 0.218 & 0.259 & 0.295 \\
UNI2 & 0.197 & 0.225 & 0.257 & 0.295 \\
H-Optimus-0 & 0.205 & 0.233 & 0.269 & 0.307 \\
\bottomrule
\end{tabular*}
\endgroup
\end{center}

\begin{itemize}
\item \textbf{Result:} every encoder improves as the training scale increases;
their mean Gene Pearson rises from 0.177 to 0.205, 0.241, and 0.278.
\end{itemize}

\textbf{A2 (Resolution---our central benchmark).}
\begin{itemize}
\item \textbf{Task and purpose:} this benchmark asks whether single-cell gene
expression can be predicted directly from histology. The input is the histology
image alone, and the output is the RNA-expression vector of the aligned cell.
Such predictions can support cell-type or cell-state assignment and other
cell-resolved downstream analyses.
\item \textbf{Evaluation protocols:} we designed in-sample prediction,
cross-sample/cross-patient prediction, and cross-platform prediction to separate
increasingly difficult forms of generalization.
\item \textbf{Definitions:} we first implemented cell-level baseline
experiments and then proposed STBoost, a model-agnostic framework that adapts
spot-level methods to cell-level prediction. To avoid ambiguity, every BLEEP
result in our experiments uses BLEEP after this STBoost adaptation. STBoost-Ref
is our further optimized reference method, which incorporates a diffusion loss
and related components detailed in the Appendix. It produces comparatively
strong results, but it also has limitations; we present it as a reference and
hope future work will develop stronger methods using this dataset.
\end{itemize}

\begin{center}
\begingroup
\footnotesize
\setlength{\tabcolsep}{4.2pt}
\renewcommand{\arraystretch}{1.04}
\textit{\textbf{Result. Per-organ cell-aligned breadth benchmark.}
Rows follow the 25-category release taxonomy. Results use top-50
training-selected HVGs, four contiguous spatial holdouts, and a 5\%
train--test buffer; the final row is the 30-sample native-Xenium macro.}\\[5pt]
\begin{tabular*}{\linewidth}{@{\extracolsep{\fill}}lrrrrrr@{}}
\toprule
\multicolumn{1}{c}{} & \multicolumn{3}{c}{Gene Pearson}
& \multicolumn{3}{c}{Image prediction} \\
\cmidrule(lr){2-4}\cmidrule(lr){5-7}
\shortstack[l]{Release category\\(evaluated species)} &
Image &
\shortstack{Coordinate\\only} &
\shortstack{Spatial\\KNN} &
\shortstack{Gene\\Spearman} &
\shortstack{Cell\\Pearson} &
F1 \\
\midrule
\textbf{Bone (mouse)} & 0.030 & \textbf{0.124} & 0.055 & 0.037 & 0.180 & 0.240 \\
Brain (human + mouse) & \textbf{0.399} & 0.010 & 0.036 & 0.371 & 0.570 & 0.752 \\
Breast (human) & \textbf{0.400} & 0.053 & 0.035 & 0.375 & 0.433 & 0.487 \\
Cervical (human) & \textbf{0.178} & 0.017 & 0.014 & 0.157 & 0.247 & 0.026 \\
\textbf{Colon (mouse)} & \textbf{0.615} & -0.137 & 0.072 & 0.588 & 0.739 & 0.744 \\
Colorectal (human) & \textbf{0.431} & 0.056 & 0.083 & 0.419 & 0.532 & 0.623 \\
\textbf{Embryo (mouse)} & \textbf{0.278} & 0.037 & 0.045 & 0.251 & 0.517 & 0.342 \\
Head (zebrafish) & \textbf{0.216} & 0.022 & 0.031 & 0.193 & 0.429 & 0.287 \\
Heart (human) & \textbf{0.292} & 0.010 & 0.021 & 0.266 & 0.688 & 0.304 \\
Kidney (human) & \textbf{0.402} & 0.019 & 0.017 & 0.361 & 0.471 & 0.460 \\
Liver (human) & -0.002 & 0.003 & \textbf{0.010} & -0.002 & 0.561 & 0.400 \\
Lung (human) & \textbf{0.359} & 0.065 & 0.053 & 0.317 & 0.402 & 0.457 \\
Lymph Node (human) & \textbf{0.141} & 0.043 & 0.013 & 0.131 & 0.164 & 0.023 \\
Ovarian (human) & \textbf{0.250} & 0.122 & 0.104 & 0.239 & 0.312 & 0.234 \\
Ovarian glands (human) & \textbf{0.402} & 0.059 & 0.043 & 0.391 & 0.559 & 0.756 \\
Pancreas (human) & \textbf{0.324} & 0.090 & 0.068 & 0.272 & 0.505 & 0.275 \\
Pancreatic (human) & \textbf{0.366} & 0.080 & 0.051 & 0.341 & 0.540 & 0.694 \\
Pancreatic duct gland (human) & \textbf{0.331} & 0.091 & 0.100 & 0.282 & 0.408 & 0.386 \\
Plant (\textit{A. thaliana}) & -0.011 & 0.004 & \textbf{0.012} & -0.009 & 0.385 & 0.052 \\
Prostate (human) & \textbf{0.263} & -0.015 & 0.025 & 0.244 & 0.263 & 0.083 \\
Seed (soybean) & -0.020 & -0.003 & \textbf{0.008} & -0.018 & 0.314 & 0.037 \\
Skin (human) & \textbf{0.359} & 0.048 & 0.086 & 0.299 & 0.437 & 0.368 \\
\textbf{Small Intestine (mouse)} & \textbf{0.572} & -0.104 & 0.067 & 0.548 & 0.701 & 0.715 \\
Tonsil (human) & \textbf{0.294} & 0.028 & 0.018 & 0.249 & 0.462 & 0.399 \\
Xenograft (human + mouse) & \textbf{0.387} & 0.041 & 0.049 & 0.362 & 0.526 & 0.589 \\
\midrule
\textbf{30-sample macro} & \textbf{0.324} & 0.046 & 0.053 &
\textbf{0.291} & \textbf{0.442} & \textbf{0.413} \\
\bottomrule
\end{tabular*}

\parbox{0.98\linewidth}{\scriptsize\textit{Among category rows, bold names are
mouse-only. Bold Gene-Pearson values mark the best of image, coordinate, and
spatial KNN. Native-Xenium top-200 image Gene Pearson is 0.202; the top-50 95\%
CI is 0.277--0.368.}}
\endgroup
\end{center}

\textbf{A3 (Rich context).}
\begin{itemize}
\item \textbf{Purpose:} our rich-context metadata---including age, sex/gender
where reported, disease, and underlying clinical conditions---enable us to test
whether a model performs consistently across biologically and clinically
meaningful subgroups. Average performance alone can hide severe subgroup
failures, so such context-aware evaluation is an important dimension for
assessing foundation models and other predictive methods.
\item \textbf{Main-text example:} the ovarian experiment is one illustrative
case. AK (60+) \(\to\) AD (60+) versus AK (60+) \(\to\) AL (\(<40\)) changes F1
from 0.289 to 0.072, gene Spearman from 0.036 to 0.006, cell Spearman from 0.206
to 0.095, and nonzero rate from 0.166 to 0.027, revealing a large age-associated
performance disparity.
\item \textbf{Broader evidence:} we conducted analogous audits for other
foundation models and contexts. The table below shows substantial assay-,
dataset-, and site-associated gaps for Geneformer, scGPT, CONCH, UNI, and
H-Optimus-0:
\end{itemize}

\begin{center}
\begingroup
\scriptsize
\setlength{\tabcolsep}{2pt}
\renewcommand{\arraystretch}{1.05}
\begin{tabularx}{\linewidth}{@{}p{0.11\linewidth}X p{0.18\linewidth}rrr@{}}
\toprule
Domain & Task / encoder & Context; split & Avg.\ BA & Worst BA & Gap \\
\midrule
Single-cell & Bone marrow / Geneformer & assay; patient-CV
& 0.938 & 0.669 & 0.327 \\
Single-cell & Bone marrow / scGPT & assay; patient-CV
& 0.962 & 0.740 & 0.258 \\
Single-cell & Bone marrow / scVI-style & assay; patient-CV
& 0.932 & 0.616 & 0.373 \\
Single-cell & Ten tissues / scGPT & dataset; patient-CV
& 0.667--0.962 & 0.004--0.892 & max 0.975 \\
Pathology & LUAD KRAS / CONCH & site; leave-one-site
& 0.499 & 0.375 & 0.542 \\
Pathology & LGG IDH / UNI & site; leave-one-site
& 0.682 & 0.464 & 0.471 \\
Pathology & LGG IDH / H-optimus0 & site; leave-one-site
& 0.747 & 0.476 & 0.524 \\
\bottomrule
\end{tabularx}
\endgroup
\end{center}

\begin{itemize}
\item \textbf{Significance and boundary:} BA is balanced accuracy, and Gap is
best-to-worst performance. The purpose of this axis is to make
context-dependent performance bias measurable, not to claim that the resource
itself removes bias. Because ovarian age is nested within sample/patient and the
panels differ, this example supports an age-associated, sample-confounded shift
rather than a causal age effect.
\end{itemize}

\questionsep
\textbf{Q3.}\quad
\textit{How are bin2cell targets constructed and validated?}

\textbf{A1 (Construction).}
\begin{itemize}
\item Appendix S2 shows the official Visium HD registration of native
2-\(\mu\)m bins to H\&E.
\item Intersecting same-frame cell polygons aggregate counts; unsupported cells
are removed and conflicts resolved deterministically.
\item We will move this pipeline and an example to the main text and label
outputs \emph{derived cell-aligned targets}, not measured single-cell RNA.
\end{itemize}

\textbf{A2 (Validation).}\enspace We audited 3,000 raw CellViT polygons per
dataset (9,000 total) under 1-\(\mu\)m registration shifts and erosion/dilation:

\begin{center}
\begingroup
\scriptsize
\setlength{\tabcolsep}{2pt}
\renewcommand{\arraystretch}{1.05}
\begin{tabularx}{\linewidth}{@{}Xrrrr@{}}
\toprule
Visium HD example & Raw polygons with canonical bins & Shift bin Jaccard &
Shift expression cosine & Shift median \( |\Delta\mathrm{UMI}| \) \\
\midrule
Human lung cancer & 49.4\% & 0.727--0.733 & 0.954--0.957 & 11.1\%--11.3\% \\
Mouse brain & 97.5\% & 0.806--0.816 & 0.936--0.939 & 6.0\%--6.1\% \\
Human pancreas & 50.0\% & 0.706--0.714 & 0.994 & 13.0\%--13.7\% \\
\bottomrule
\end{tabularx}
\endgroup
\end{center}

\begin{itemize}
\item Canonical-bin percentages use all raw polygons; unsupported polygons are
excluded before release.
\item Shifts preserve expression direction but alter membership/UMIs;
erosion/dilation yields Jaccard 0.462--0.720, cosine 0.871--0.993, and
\(|\Delta\mathrm{UMI}|\) 32.8\%--66.6\%.
\item We will release parameters and per-sample QC, avoid calling aggregation
ground truth, and report native-Xenium evidence separately.
\end{itemize}

\questionsep
\textbf{Q4.}\quad
\textit{How are LMM-generated texts produced and validated?}

\textbf{A1 (Generation).}\enspace GPT-4o verbalizes supplied fields without adding attributes. We regret omitting
the completed audit while shortening this section and will restore prompts,
field checks, examples, and matrices:

\begin{center}
\begingroup
\scriptsize
\setlength{\tabcolsep}{2pt}
\renewcommand{\arraystretch}{1.08}
\begin{tabular*}{\linewidth}{@{\extracolsep{\fill}}lrrrrrr@{}}
\toprule
Model & Match \(n\) & Match sim. & Mismatch sim. & Enrichment &
Pairwise AUC & Top-1 / Top-3 \\
\midrule
PLIP & 15 & 0.262 & 0.053 & \(4.96{\times}\) & 0.993 & 84.6\% / 100\% \\
CONCH & 29 & 0.040 & 0.0021 & \(19.1{\times}\) & 0.988 & 90.9\% / 100\% \\
\bottomrule
\end{tabular*}
\endgroup
\end{center}

\textbf{A2 (Validation and scope).}
\begin{itemize}
\item Both VLMs reach AUC \(\geq 0.988\) and 100\% Top-3 retrieval.
\item These metrics assess image relevance, while source-field checks assess
factuality.
\item Text remains optional context, not ground truth.
\end{itemize}

\end{document}
